\subsection{Affine Independence and Barycentric Coordinates}

Throughout this subsection, $\mathcal{V}$ is finite-dimensional over a fixed field $\mathbb{F}$.

\begin{definition} \label{def:affinely-independent-set} \hlt{Affinely Independent Sets}\\
A nonempty set $E\subseteq\mathcal{V}$ is \hlt{affinely independent} iff $(\forall\,\mathbf{x}\in E)\,(E-\mathbf{x})\setminus\{\mathbf{0}\}$ is linearly independent.
\end{definition}

\begin{theorem} \label{thm:affine-independence-equivalences} \hlt{Affine Independence Criteria}\\
Let $\emptyset\neq E\subseteq\mathcal{V}$. For $F\subseteq\mathcal{V}$ and $n\in\mathbb{Z}_{>0}$, put $F_{\neq}^n\defeq\{(\mathbf{x}_1,\ldots,\mathbf{x}_n)\in F^n:\mathbf{x}_j\neq\mathbf{x}_k\ (j\neq k)\}$. The following conditions are equivalent.
\begin{enumerate}[label=(\roman*)]
\setlength{\itemsep}{0pt}
\item $E$ is affinely independent.
\item $(\forall\,\mathbf{x}\in E)\,\mathbf{x}\notin\operatorname{affhull}(E\setminus\{\mathbf{x}\})$, where $\operatorname{affhull}(\emptyset)=\emptyset$.
\item $(\forall\,n\in\mathbb{Z}_{>0})(\forall\,(\mathbf{x}_1,\ldots,\mathbf{x}_n)\in E_{\neq}^n)(\forall\,\mathbf{r}\in\mathbb{F}^n)\,[\sum\limits_{j=1}^n r_j\mathbf{x}_j=\mathbf{0}\land\sum\limits_{j=1}^n r_j=0\Rightarrow\mathbf{r}=\mathbf{0}]$.
\item $(\forall\,n\in\mathbb{Z}_{>0})(\forall\,(\mathbf{x}_1,\ldots,\mathbf{x}_n)\in E_{\neq}^n)(\forall\,\mathbf{r},\mathbf{s}\in\mathbb{F}^n)\,[\sum\limits_{j=1}^n r_j\mathbf{x}_j=\sum\limits_{j=1}^n s_j\mathbf{x}_j\land\sum\limits_{j=1}^n r_j=\sum\limits_{j=1}^n s_j=1\Rightarrow\mathbf{r}=\mathbf{s}]$.
\end{enumerate}
\begin{proof}
\begin{enumerate}[label=(\roman*)]
\setlength{\itemsep}{0pt}
\item[$(i)$]$\Rightarrow (ii)$ Fix $\mathbf{x}\in E$. The claim holds by the stated convention when $E=\{\mathbf{x}\}$. Otherwise, suppose $\mathbf{x}\in\operatorname{affhull}(E\setminus\{\mathbf{x}\})$.\\
By Theorem \ref{thm:affine-hull-span-formula}(i), collect equal points, $(\exists\,m\in\mathbb{Z}_{>0})(\exists\,(\mathbf{y}_1,\ldots,\mathbf{y}_m)\in(E\setminus\{\mathbf{x}\})_{\neq}^m)(\exists\,\mathbf{r}\in\mathbb{F}^m)\,[\mathbf{x}=\sum\limits_{j=1}^m r_j\mathbf{y}_j\land\sum\limits_{j=1}^m r_j=1]$.\\
Thus $\sum\limits_{j=1}^m r_j(\mathbf{y}_j-\mathbf{x})=\mathbf{0}$ with $\mathbf{r}\neq\mathbf{0}$, contrary to (i). 
\item[$(ii)$]$\Rightarrow (iii)$ Fix $n\in\mathbb{Z}_{>0}$, $(\mathbf{x}_1,\ldots,\mathbf{x}_n)\in E_{\neq}^n$, and $\mathbf{r}\in\mathbb{F}^n$ with $\sum\limits_{j=1}^n r_j\mathbf{x}_j=\mathbf{0}$ and $\sum\limits_{j=1}^n r_j=0$. If $\mathbf{r}\neq\mathbf{0}$, choose $k$ with $r_k\neq0$.\\
Then $n\geq2$, $\mathbf{x}_k=\sum\limits_{j\neq k}(-\frac{r_j}{r_k})\mathbf{x}_j$, and $\sum\limits_{j\neq k}-\frac{r_j}{r_k}=1$, so $\mathbf{x}_k\in\operatorname{affhull}(E\setminus\{\mathbf{x}_k\})$ by Theorem \ref{thm:affine-hull-span-formula}(i), contrary to (ii). Hence $\mathbf{r}=\mathbf{0}$. 
\item[$(iii)$]$\Rightarrow (iv)$ After subtraction, $\sum\limits_{j=1}^n(r_j-s_j)\mathbf{x}_j=\mathbf{0}$ and $\sum\limits_{j=1}^n(r_j-s_j)=0$, so $\mathbf{r}=\mathbf{s}$ by (iii). \item[$(iv)$]$\Rightarrow (i)$ Fix $\mathbf{x}\in E$. The empty list is linearly independent. Fix $m\in\mathbb{Z}_{>0}$, $(\mathbf{y}_1,\ldots,\mathbf{y}_m)\in(E\setminus\{\mathbf{x}\})_{\neq}^m$, and $\mathbf{a}\in\mathbb{F}^m$ with $\sum\limits_{j=1}^m a_j(\mathbf{y}_j-\mathbf{x})=\mathbf{0}$. On $(\mathbf{y}_1,\ldots,\mathbf{y}_m,\mathbf{x})\in E_{\neq}^{m+1}$, put $\mathbf{c}\defeq(a_1,\ldots,a_m,-\sum\limits_{j=1}^m a_j)$ and $\mathbf{b}\defeq(0,\ldots,0,1)$.\\
Both $\mathbf{b}$ and $\mathbf{b}+\mathbf{c}$ have coefficient sum $1$ and represent $\mathbf{x}$, so $\mathbf{c}=\mathbf{0}$ by (iv). Hence $\mathbf{a}=\mathbf{0}$, and $(E-\mathbf{x})\setminus\{\mathbf{0}\}$ is linearly independent.
\end{enumerate}
\hlt{Structural Commentary}: Cycle through excluded hulls, zero-sum relations, coordinate uniqueness, and translated independence.\\
\hlt{Key Insight}: Affine dependence is exactly a nonzero coefficient relation with total sum zero.
\end{proof}
\end{theorem}

\begin{corollary} \label{cor:finite-affine-independence-dimension-criterion} \hlt{Dimension Criterion for Affine Independence}\\
Let $E\subseteq\mathcal{V}$ be finite and nonempty with $\abs{E}=n$. Then $E$ is affinely independent iff $\dim(\operatorname{affhull}(E))=n-1$.
\begin{proof}
\begin{enumerate}[label=\roman*.]
\setlength{\itemsep}{0pt}
\item[$(\Rightarrow)$] Fix $\mathbf{x}\in E$ and put $B_{\mathbf{x}}\defeq(E-\mathbf{x})\setminus\{\mathbf{0}\}$. The set $B_{\mathbf{x}}$ has $n-1$ elements and is linearly independent by Definition \ref{def:affinely-independent-set}; hence any ordering is a basis of $\operatorname{span}(E-\mathbf{x})$. Thus $\dim(\operatorname{affhull}(E))=\dim(\operatorname{span}(E-\mathbf{x}))=n-1$ by Theorem \ref{thm:affine-hull-span-formula}(ii) and Definition \ref{def:dimension}.
\item[$(\Leftarrow)$] Fix $\mathbf{x}\in E$ and put $B_{\mathbf{x}}\defeq(E-\mathbf{x})\setminus\{\mathbf{0}\}$. By Theorem \ref{thm:affine-hull-span-formula}(ii), $\dim(\operatorname{span}(E-\mathbf{x}))=n-1$.\\
Any ordering of the $n-1$ elements of $B_{\mathbf{x}}$ spans this space, so it is a basis by Proposition \ref{prop:full-dimension-basis-criteria}(i).\\
Thus $B_{\mathbf{x}}$ is linearly independent, and $E$ is affinely independent by Definition \ref{def:affinely-independent-set}.
\end{enumerate}
\hlt{Structural Commentary}: Translate one point to zero, then compare the remaining spanning list with the hull dimension.\\
\hlt{Key Insight}: Removing one affine origin leaves exactly $n-1$ direction vectors.
\end{proof}
\end{corollary}

\begin{definition} \label{def:affine-basis-barycentric-coordinates} \hlt{Affine Bases and Barycentric Coordinates}\\
Let $A=\mathbf{x}+\mathcal{S}$ be an $n$-flat.
\begin{enumerate}[label=(\roman*)]
\setlength{\itemsep}{0pt}
\item A finite subset $\mathcal{B}\subseteq A$ is an \hlt{affine basis} of $A$ iff $\mathcal{B}$ is affinely independent and $\operatorname{affspan}(\mathcal{B})=A$.
\item A tuple $(\mathbf{p}_1,\ldots,\mathbf{p}_m)$ is an \hlt{ordered affine basis} of $A$ iff its entries are distinct and $\{\mathbf{p}_1,\ldots,\mathbf{p}_m\}$ is an affine basis of $A$.
\item If $(\mathbf{p}_1,\ldots,\mathbf{p}_m)$ is an ordered affine basis of $A$ and $\mathbf{v}\in A$,\\
then $\mathbf{r}=(r_1,\ldots,r_m)\in\mathbb{F}^m$ is a \hlt{barycentric coordinate vector} of $\mathbf{v}$ iff $\mathbf{v}=\sum\limits_{j=1}^{m}r_j\mathbf{p}_j$ and $\sum\limits_{j=1}^{m}r_j=1$.
\end{enumerate}
\end{definition}

\begin{theorem} \label{thm:affine-bases-and-barycentric-coordinate-uniqueness} \hlt{Affine Bases and Barycentric Coordinates}\\
Let $A=\mathbf{x}+\mathcal{S}$ be an $n$-flat.
\begin{enumerate}[label=(\roman*)]
\setlength{\itemsep}{0pt}
\item If $(\mathbf{s}_1,\ldots,\mathbf{s}_n)$ is an ordered basis of $\mathcal{S}$, then $(\mathbf{x}+\mathbf{s}_1,\ldots,\mathbf{x}+\mathbf{s}_n,\mathbf{x})$ is an ordered affine basis of $A$.
\item If $\mathcal{B}$ is an affine basis of $A$, then $\abs{\mathcal{B}}=n+1$.
\item If $(\mathbf{p}_1,\ldots,\mathbf{p}_{n+1})$ is an ordered affine basis of $A$, then $(\forall\,\mathbf{v}\in A)(\exists!\,\mathbf{r}\in\mathbb{F}^{n+1})\,[\mathbf{v}=\sum\limits_{j=1}^{n+1}r_j\mathbf{p}_j\land\sum\limits_{j=1}^{n+1}r_j=1]$.
\end{enumerate}
\begin{proof}
\begin{enumerate}[label=(\roman*)]
\setlength{\itemsep}{0pt}
\item Put $\mathbf{p}_j\defeq\mathbf{x}+\mathbf{s}_j$ for $1\leq j\leq n$ and $\mathbf{p}_{n+1}\defeq\mathbf{x}$. Every zero-sum relation on a subtuple of $(\mathbf{p}_1,\ldots,\mathbf{p}_{n+1})$ extends by zero coefficients to the full tuple. If $\mathbf{a}\in\mathbb{F}^{n+1}$ satisfies $\sum\limits_{j=1}^{n+1}a_j\mathbf{p}_j=\mathbf{0}$ and $\sum\limits_{j=1}^{n+1}a_j=0$, then $\sum\limits_{j=1}^{n}a_j\mathbf{s}_j=\mathbf{0}$.\\
Hence $a_1=\cdots=a_n=0$ by independence of the $\mathbf{s}_j$, and $a_{n+1}=0$ by the coefficient sum.\\
Thus $\{\mathbf{p}_1,\ldots,\mathbf{p}_{n+1}\}$ is affinely independent by Theorem \ref{thm:affine-independence-equivalences}(iii).\\
The identity $\operatorname{affspan}(\{\mathbf{p}_1,\ldots,\mathbf{p}_{n+1}\})=\mathbf{x}+\operatorname{span}(\mathbf{s}_1,\ldots,\mathbf{s}_n)=A$ holds by Theorem \ref{thm:affine-hull-span-formula}(i); tuple is an ordered affine basis.
\item Put $m\defeq\abs{\mathcal{B}}$. Then $\dim(\operatorname{affhull}(\mathcal{B}))=m-1$ by Corollary \ref{cor:finite-affine-independence-dimension-criterion}.\\
By Definition \ref{def:affine-basis-barycentric-coordinates}(i) and Theorem \ref{thm:affine-hull-span-formula}(i), $\operatorname{affhull}(\mathcal{B})=\operatorname{affspan}(\mathcal{B})=A$. Thus $n=m-1$, so $m=n+1$.
\item Fix $\mathbf{v}\in A$. Since $\operatorname{affspan}(\{\mathbf{p}_1,\ldots,\mathbf{p}_{n+1}\})=A$, an affine combination of these basis points represents $\mathbf{v}$ by Definition \ref{def:affine-hull-and-affine-span}(ii).\\
Combining repeated coefficients and assigning zero to omitted points produces $\mathbf{r}\in\mathbb{F}^{n+1}$ with $\mathbf{v}=\sum\limits_{j=1}^{n+1}r_j\mathbf{p}_j$ and $\sum\limits_{j=1}^{n+1}r_j=1$.\\
This vector is unique by Theorem \ref{thm:affine-independence-equivalences}(iv).
\end{enumerate}
\hlt{Structural Commentary}: Adjoin an affine origin to a linear basis, compare dimensions, then invoke affine-representation uniqueness.\\
\hlt{Key Insight}: One extra basis point records the missing affine origin.
\end{proof}
\end{theorem}

\begin{example} \label{ex:barycentric-coordinates-on-plane} \hlt{Barycentric Coordinates on a Plane}\\
Let $A=\mathbf{x}+\operatorname{span}(\mathbf{v}_1,\mathbf{v}_2)\subseteq\mathbb{R}^3$, where $(\mathbf{v}_1,\mathbf{v}_2)$ is linearly independent.\\
Put $\mathbf{p}_1\defeq\mathbf{x}+\mathbf{v}_1$, $\mathbf{p}_2\defeq\mathbf{x}+\mathbf{v}_2$, and $\mathbf{p}_3\defeq\mathbf{x}$. Then $(\mathbf{p}_1,\mathbf{p}_2,\mathbf{p}_3)$ is an ordered affine basis of $A$ by Theorem \ref{thm:affine-bases-and-barycentric-coordinate-uniqueness}(i).\\
If $\mathbf{a}=\mathbf{x}+t_1\mathbf{v}_1+t_2\mathbf{v}_2$, then $\mathbf{a}=t_1\mathbf{p}_1+t_2\mathbf{p}_2+(1-t_1-t_2)\mathbf{p}_3$ and $t_1+t_2+(1-t_1-t_2)=1$.\\
Hence the barycentric coordinate vector of $\mathbf{a}$ is $(t_1,t_2,1-t_1-t_2)$ by Theorem \ref{thm:affine-bases-and-barycentric-coordinate-uniqueness}(iii).
\end{example}
